% ============================================================
%  MÓDULO — Infraestrutura, Serviços e Publicação
%  Escrito na revisão R616 (substitui o stub "Em construção").
% ============================================================
\chapter{Infraestrutura, Serviços e Publicação}

\roteirodidatico
{Um programa depende do ambiente, dos serviços a que se conecta e de regras que controlam operações. Este capítulo mostra como reconhecer essas dependências na prática.}
{\emph{Serviço} executa uma função acessível por interface; \emph{deploy} publica uma versão; \emph{sondagem} consulta um serviço; \emph{hook} intercepta um evento; \emph{perfil} seleciona dependências.}
{Separe instalação local, serviço externo e publicação; em cada etapa confira configuração, versão, credenciais, resposta e registro associado.}
{Uma checagem de saúde confirma uma resposta limitada naquele instante. Ela não certifica segurança, disponibilidade futura, integridade de ponta a ponta ou qualidade do produto.}
{Que parte da cadeia foi efetivamente exercitada e quais riscos ou condições ficaram fora da verificação?}

\casoguiado{da instalação à publicação}
{Uma compilação local bem-sucedida mostra que o documento foi produzido naquele ambiente. Uma publicação exige passos adicionais: enviar o artefato, aguardar o serviço e consultar a URL resultante. Uma resposta HTTP confirma uma condição de acesso naquele instante; não prova que todos os leitores verão a versão correta nem que a publicação continuará disponível. Registre artefato, destino, horário e resposta observada.}

\elemento{Infraestrutura do Core}{\metainfo{ID}{I-00} \hfill
\metainfo{Camada}{8 — Infraestrutura} \hfill
\metainfo{Rota}{\pth{integrations/}, \pth{.opencode/}}}

\begin{descricao}
Os módulos anteriores descrevem o \emph{pensar} do Core: orquestração, memória,
qualidade, raciocínio, ciência. Este descreve o \emph{operar} --- o que existe
\emph{fora} do código do orquestrador e o faz funcionar no mundo real: o ambiente,
os serviços que ficam de pé, as guardas que interrompem operações perigosas e a
cadeia que leva um artefato do repositório até um site publicado.
\end{descricao}

\begin{conceito}
A infraestrutura organiza-se em três camadas, da mais estável para a mais externa:
\begin{enumerate}[leftmargin=1.4em,itemsep=1pt]
  \item \textbf{Ambiente} --- dependências por perfil e instalação reprodutível
        (desenvolvida no Módulo de Instalação).
  \item \textbf{Serviços} --- servidores MCP que expõem capacidades a qualquer
        cliente compatível.
  \item \textbf{Guardas e publicação} --- hooks que interrompem, plugins que
        estendem e a cadeia de deploy com sondagem real.
\end{enumerate}
\end{conceito}

\begin{funcao}
Transformar ``funciona na minha máquina'' em ``funciona em qualquer máquina
auditável'' --- a mesma regra do capítulo de instalação, agora aplicada a serviços
e à publicação.
\end{funcao}

\begin{niveis}
\textbf{Intuição:} a parte invisível que faz a peça acontecer --- tomada, motor e
freio de mão.\\
\textbf{Implementação:} venv por perfil, 7 servidores MCP, 4 hooks, 2 plugins e cadeia
de deploy com sondagem HTTP real.\\
\textbf{Formalização:} infraestrutura auditável, com proveniência por artefato e guardas
que impõem política sem depender da disciplina do operador.
\end{niveis}

\section{O ambiente: perfis e instalação reproduzível}

\begin{descricao}
O \pth{Makefile} oferece perfis de dependência: \pth{make install} (ambiente
virtual), \pth{make setup} (base), \pth{make setup-dev} (desenvolvimento) e
\pth{make setup-science} (perfil científico). Perfis separados existem para que a
dependência de uma tarefa científica não contamine o ambiente de operação --- o
equivalente informacional de separar o reagente de teste do reagente em uso.
\end{descricao}

\begin{metodo}
Ritual de entrada: (1) \pth{make install}; (2) \pth{make setup}; (3)
\pth{python3 -m marceloclaro.cli doctor}; (4) só então operar. O \pth{doctor} é o
exame de saúde do aparato: ambiente verde significa ``pronto para operar'', nunca
``resultado verificado''.
\end{metodo}

\section{Serviços: os sete servidores MCP}

\begin{resultado}
Servidores MCP declarados em \pth{opencode.json} (chave \pth{mcp}):
\par\smallskip
{\footnotesize
\begin{tabularx}{\textwidth}{lY}
\toprule
Servidor & Papel \\
\midrule
\pth{metacognitive-interconnect} & Blackboard, MetaBus e registro de agentes \\
\pth{litert-lm} & modelos on-device (Gemma 4 / Qwen3) via LiteRT-LM \\
\pth{antigravity-bridge} & delegação a agentes externos, busca e leitura de URL \\
\pth{scanners-mcp} & auditoria (rigor científico, literário, integridade) \\
\pth{pypi-search} & busca e avaliação de pacotes Python \\
\pth{colibri-mcp} & geração local com o motor Colibri \\
\pth{web-deploy-mcp} & sondagem de URLs, feeds e orçamento de peso \\
\bottomrule
\end{tabularx}\par}
\end{resultado}

\begin{metodo}
Política de invocação: um servidor ausente \emph{não} derruba o Core. As skills
que dependem de executor externo validam presença e versão antes do uso e degradam
para um aviso explícito --- a invocação tolerante descrita no Módulo de Integrações.
\end{metodo}

\begin{aviso}
Disponibilidade declarada $\neq$ disponibilidade real: um servidor configurado pode
estar sem credencial ou sem binário. A verificação antes do uso é obrigatória.
\end{aviso}

\section{Guardas: hooks e plugins}

\begin{resultado}
Hooks em \pth{.opencode/hooks/} e plugins em \pth{.opencode/plugins/}:
\par\smallskip
{\footnotesize
\begin{tabularx}{\textwidth}{lY}
\toprule
Componente & Função \\
\midrule
\pth{budget_guard.sh} & interrompe a operação quando o orçamento de tokens é excedido \\
\pth{credential_guard.sh} & bloqueia operação sem credencial configurada \\
\pth{js_smoke_dom.sh} (+ \cod{smoke_dom_stub.js}) & executa o JavaScript de renderização em Node antes de declarar deploy pronto \\
\pth{deploy-guards.ts} (plugin) & guardas de deploy: peso, remoção (404) e integridade do feed \\
\pth{litert-lm-provider.ts} (plugin) & registra o provedor LiteRT-LM no cliente \\
\bottomrule
\end{tabularx}\par}
\end{resultado}

\begin{conceito}
Um hook é uma \textbf{política executada antes do ato}. Sua vantagem sobre a
documentação é simples e mensurável: a documentação depende de o operador lembrar;
o hook não depende.
\end{conceito}

\section{Publicação com rastro: a cadeia de deploy}

\begin{descricao}
Publicar é a operação em que a maioria dos erros é \emph{invisível}: o comando
``passa'' e o site está quebrado. Por isso a cadeia de deploy do Core não aceita
``publicado'' como prova --- ela \emph{sonda}.
\end{descricao}

\begin{metodo}
\begin{enumerate}[leftmargin=1.4em,itemsep=1pt]
  \item \textbf{Status do build} --- consulta à API do GitHub Pages
        (\pth{web_deploy_pages_status}).
  \item \textbf{Sondagem} --- requisição \emph{GET} com \emph{Range} em cada URL,
        buscando um marcador de conteúdo (\pth{web_deploy_probe_url}).
  \item \textbf{Verificação de remoção} --- URLs retiradas devem responder 404
        (\pth{web_deploy_assert_gone}); qualquer status diferente de 404 é falha.
  \item \textbf{Orçamento de peso} --- limite de peso do diretório e de arquivo
        (\pth{web_deploy_site_weight}).
  \item \textbf{Feed} --- validação de RSS/Atom, com conferência do enclosure
        (\pth{web_deploy_validate_feed}).
\end{enumerate}
\end{metodo}

\begin{aviso}
\textbf{Lição de infraestrutura (R580, reafirmada em R616).} A sondagem usa
\emph{GET} com \cod{Range} e \emph{nunca} \cod{HEAD}: o GitHub Pages responde de
modo inconsistente a \cod{HEAD} para arquivos binários, produzindo falso 404.
Da mesma natureza é o smoke test de JavaScript com stub de DOM: \cod{node --check}
valida sintaxe e \emph{não} executa o DOM, de modo que erros de renderização só
aparecem em tempo de execução.
\end{aviso}

% ============================================================

%  Gerado por gen_mod14_ativacao.py (R613) — seção de ativação e cálculo
%  para o módulo mod-09-infra.tex. Edições manuais são preservadas nas próximas
%  execuções (marcador impede reescrita).
% ============================================================

% ATIVACAO-R613
\section[Ativação e cálculo]{Ativação e cálculo — Infraestrutura}
\elemento{ativ-09i}{\metainfo{Ciclo}{R613} \hfill \metainfo{Módulo}{mod-09-infra.tex}}

\begin{processo}
GATILHO. \pth{doctor}, serviços MCP, integrações de dados.
ROTA. serviços determinísticos; gate: rastro e proveniência por operação.
FUNÇÃO. executar serviços de suporte com rastro completo e disponibilidade mensurável.
\end{processo}

\begin{resultado}
CÁLCULO. disponibilidade = operações com rastro / operações totais; proveniência = metadados de origem/versão por artefato.
\end{resultado}

\begin{descricao}
JUSTIFICATIVA TEÓRICA. a proveniência detalhada e o rastro por operação são o que torna a infraestrutura auditável. O lastro teórico vem da citação que segue: recorte conferido em 29/09/2026, com a página de origem e tradução livre e fiel.
\end{descricao}

\begin{aviso}
LIMITE E ANTI-OVERCLAIM. rastro documenta, não garante correção do mundo externo. A verificação atesta a existência e a
localização da fonte (R110), não o mérito da afirmação.
\end{aviso}

\noindent\textit{Interpretação:} Os princípios FAIR transformam a proveniência em requisito operacional: o artefato deve ser localizável, acessível sob condições declaradas, interoperável e reutilizável. Como diagnóstico interno — não como certificação FAIR — pode-se medir $C_m=n_p/n_r$, em que $n_p$ é o número de campos obrigatórios preenchidos e $n_r$ o total requerido. $C_m=1$ indica completude do esquema adotado, não correção do conteúdo; esta ainda depende da conferência da fonte e da validação do resultado \citref{WILKINSON, M. D. et al. The FAIR Guiding Principles for scientific data management and stewardship. \emph{Scientific Data}, v.~3, 160018, 2016}{10.1038/sdata.2016.18}{p.~160018 (resumo, p.~160018). é o padrão de auditabilidade dos serviços de infraestrutura: todo dado usado pelo Core deve ser localizável, acessível, interoperável e reutilizável — o que o rastro de operações materializa.}{Good data management is not a goal in itself, but is the key conduit leading to knowledge discovery and innovation}{O bom gerenciamento de dados não é um fim em si mesmo, mas o principal canal que conduz à descoberta de conhecimento e à inovação.}{WILKINSON et al., 2016, p.~160018}\section[Leitura guiada]{Leitura guiada — Infraestrutura: recursos e restrições}

\elemento{Leitura guiada — Infraestrutura}{\metainfo{ID}{N-09} \hfill \metainfo{Ciclo}{R616} \hfill \metainfo{Estilo}{medir o que existe fora do código}}

\begin{descricao}
\textbf{O fenômeno.} Escreva no papel a frase que todo desenvolvedor já disse:
``na minha máquina funciona''. Ela é a antítese da reprodutibilidade. Este
capítulo trata do oposto: tornar \emph{verificável} aquilo que vive fora do
próprio código --- serviços, guardas e publicação.
\end{descricao}

\begin{definicao}
\textbf{Grandezas e definições.}
\begin{itemize}[leftmargin=1.3em,itemsep=1pt]
  \item \textbf{Disponibilidade} ($d$): operações com rastro $/$ operações totais.
  \item \textbf{Proveniência} ($p$): metadados de origem e versão por artefato.
  \item \textbf{Orçamento de peso} ($w$): limite de peso do diretório ($1\,000$ MB) e de arquivo ($25$ MB).
  \item \textbf{Marcador de conteúdo} ($m$): trecho esperado na resposta HTTP, que distingue ``200 com a página'' de ``200 com a página de erro''.
\end{itemize}
\end{definicao}

\begin{metodo}
\textbf{Dedução passo a passo.} Por que sondar em vez de acreditar no comando de
publicação?
\begin{enumerate}[leftmargin=1.4em,itemsep=2pt]
  \item ``Publicado'' é uma afirmação do cliente; a prova é a resposta do servidor.
  \item Um \cod{HEAD} devolve só cabeçalho e, no GitHub Pages, mente para binários --- falso 404.
  \item Um \cod{GET} com \cod{Range} lê de fato o início do arquivo e custa cerca de 1 KB.
  \item Comparar o corpo com o marcador $m$ separa página servida de página de erro.
  \item A remoção de um artefato só está concluída quando a URL responde 404; status outro que 404 é falha por definição do gate.
\end{enumerate}
\end{metodo}

\begin{resultado}
\textbf{Exemplo resolvido (guarda de peso).} Considere um pacote de publicação com
$68$ arquivos, $1\,240$ MB no total e um único arquivo de áudio de $1\,180$ MB:
\begin{enumerate}[leftmargin=1.4em,itemsep=1pt]
  \item Guarda de total: $1\,240 > 1\,000$ MB $\Rightarrow$ \textbf{denegado} (limite de diretório).
  \item Guarda de arquivo: $1\,180 > 25$ MB $\Rightarrow$ \textbf{denegado} (limite de arquivo).
  \item Correção: dividir ou reprocessar o áudio. A segunda falha é a mais dura, porque obriga a mudar a arquitetura do feed, e não apenas um parâmetro de build.
\end{enumerate}
\textbf{Ordem de grandeza:} as guardas atuam em duas escalas --- megabytes
(site) e bytes (conteúdo) --- enquanto a sonda \cod{GET}+\cod{Range} custa
$\approx 1$ KB por URL, cerca de mil vezes menos que um download integral.
\end{resultado}

\begin{resultado}
\textbf{Ordem de grandeza e verificação.} A publicação tem um orçamento explícito: o limite padrão é $1000$ MB por diretório e $25$ MB por arquivo, o que é da ordem de grandeza de um repositório de sites estáticos --- mas o orçamento é a \emph{folga}, não o consumo medido. Verificação: \cod{python3 -c "import pathlib; print(sum(f.stat().st_size for f in pathlib.Path('.').rglob('*') if f.is_file())/1e6)"} em MB. O ponto didático é a proporção: entre o orçamento e o consumo há uma razão de folga típica de $1{:}10$ a $1{:}100$, e é essa folga que permite publicar sem reescrever. Um sistema que consome mais de metade do orçamento deixa de ser publicável com segurança --- e o PDF do próprio livro, com suas centenas de páginas, é o exemplo do artefato que mais pesa e o menos se beneficia de ser comprimido.
\end{resultado}

\begin{aviso}
\textbf{Limite de validade.} Os valores de $1\,240$ MB e $1\,180$ MB são
\emph{ilustrativos}. O gate verifica presença, marcador e peso; ele não verifica
se o conteúdo publicado está \emph{correto} editorialmente.
\end{aviso}

\begin{resultado}
\textbf{Exercícios.} (1) Por que \cod{HEAD} é proibido na sondagem? (2) Um
arquivo de $24$ MB passa na guarda de arquivo: passa na de total? (3) Que
evidência extra um feed RSS precisaria para ser considerado íntegro?
\end{resultado}

\section{Resumo e exercícios}

\begin{resultado}
\textbf{O que você deve ter aprendido:} (1) infraestrutura é o ``operar'', distinto
do ``pensar''; (2) perfis de dependência separam ambientes; (3) 7 servidores MCP,
4 hooks e 2 plugins formam a camada executável; (4) publicar exige sondagem
(\cod{GET}+\cod{Range}), nunca \cod{HEAD}; (5) removido só está removido quando
a URL responde 404. \textbf{Exercício sugerido:} rode o \pth{doctor}, anote as
verificações e escolha uma delas para reproduzir manualmente (por exemplo, contar
os agentes registrados no Blackboard).
\end{resultado}
